% 本文件由 LaTeX 排版 Agent 根据有效 translation.json 自动生成。
% author=John Chrysostom; work_id=1915; title=论欧特罗皮乌斯讲道词 第二篇

\chapter{论\Person{欧特罗皮乌斯}讲道词 第二篇}

% structure: p0001 source=h1 target=omitted reason=page-title-already-rendered-by-parent

\Emph{\Person{欧特罗皮乌斯}在教会之外被发现并被俘之后。}\par

1. 草地和花园确实令人愉悦，但研读圣经更令人愉悦。因为在那里，花朵会凋谢，而在这里，思想常开不败；那里是和风，这里是圣神的嘘气；那里是荆棘树篱，这里却是天主的守护眷顾；那里是蝉鸣，这里是先知们的旋律；那里是视觉的快乐，这里则是研读的益处。花园局限于一个地方，而圣经则遍及全世界；花园受季节变迁的影响，而圣经在冬夏都枝叶茂盛、果实累累。让我们尽心致力于研读圣经吧：因为如果你这样做，圣经将驱散你的沮丧，生发愉悦，根除恶习，使美德扎根，并且在生活的喧嚣中，它将拯救你免于像被汹涌波浪抛掷的人那样受苦。大海咆哮，但你却航向平静；因为你以研读圣经为你的领航员；因为这是一根生活的烦恼不能折断的缆绳。我现在所说并非谎言，事件本身便是见证。几天前，教会曾被围困：一支军队来了，眼中喷火，却未能烧焦橄榄树；刀剑出鞘，却无人受伤；皇门危殆，但教会安然无恙。然而战争的浪潮涌向这里；因为这里寻找避难者，我们抵挡了他们，不畏惧他们的愤怒。这是为什么呢？因为我们持有确实的保证，即这句话：\BibleQuote{你是伯多禄，在这磐石上，我要建立我的教会：阴间的门决不能战胜她。}\BibleRef{玛 16:18}当我说教会时，我指的不仅是一个地方，也是一种生活方式：我不是指教会的墙壁，而是指教会的法律。当你投奔教会时，不要仅仅寻求地方的庇护，而是寻求地方的精神。因为教会不是墙壁和屋顶，而是信德和生命。\par

不要告诉我，这个人被交出是因为教会将他交出；如果他未抛弃教会，他就不会被交出。不要说，他逃到这里避难，然后被交出：教会没有抛弃他，而是他抛弃了教会。他不是从教会内部被交出，而是在教会墙外。他为何离弃了教会？你想要拯救自己吗？你应当紧紧抓住祭台。这里没有墙壁，但有天主的守护眷顾。你是个罪人吗？天主不会拒绝你：因为\BibleQuote{他不是来召义人，而是来召罪人悔改。}\BibleRef{玛 9:13}那个娼妓得救了，因为她紧抱他的脚。你们今天听了这个诵读吗？我说这些，是要你们不要犹豫去投奔教会。与教会同在，教会不会把你交给敌人；但如果你从教会逃走，教会不是被捕的原因。因为如果你在羊栈内，狼就不会进入；但如果你出去，你就会成为野兽的猎物：然而这不是羊栈的过错，而是你自己的怯懦。教会没有脚。不要和我谈墙壁和武器：因为墙壁随着时间旧了，但教会永不衰老。墙壁被野蛮人击碎，但连魔鬼也不能胜过教会。我的话并非虚张声势，事实就是证据。多少人攻击过教会，然而攻击者都灭亡了，而教会本身却高升到天上？教会具有这样的力量：当她被攻击时，她得胜；当为她设下网罗时，她得胜；当她受侮辱时，她的繁荣增长；她受伤却不因伤沉没；被波浪抛掷却不被淹没；被风暴袭击却未沉船；她角力却不被击败，争战却不被征服。那么，她为什么允许这场战争发生呢？是为了更显明她胜利的光荣。那一天你在场，你看见了什么武器被调动来反对她，士兵的愤怒比火更猛烈，我被迫赶往皇宫。但这又有什么呢？因天主的恩宠，那些事没有一件使我惊慌。\par

2. 我现在说这些，是要你们也效法我的榜样。但我为何不沮丧呢？因为我不惧怕任何眼前的恐怖。有什么可怕的呢？死亡？不，这并不可怕：因为我们很快便抵达风平浪静的港湾。或是被夺去财物？\BibleQuote{我赤身出于母胎，也必赤身归去；}\BibleRef{约 1:21}或是流放？\BibleQuote{大地属于主，及其中的万物；}\BibleRef{玛 5:12}或是诬告？\BibleQuote{你们欢喜踊跃罢，因为几时人为了我而辱骂迫害你们，捏造一切坏话毁谤你们，你们在天上的赏报是丰厚的。}\BibleRef{玛 5:12}我看到刀剑，就默想天堂；我预料死亡，就想起复活；我注视此世的痛苦，就思量天上的奖赏；我观察敌人的阴谋，就默想天上的冠冕：因为争战的时刻足以激励和安慰。诚然！我被强行拖走，但并未因此受辱，因为真正的羞辱只有一种，就是罪：即使全世界侮辱你，若你未侮辱自己，你就未受侮辱。真正的背叛是良心的背叛：不要背叛你的良心，就没有人能背叛你。我被拖走，我看到这些事件——更确切地说，我看到我的话语化成了事件，我看到我在言语中宣讲的言论，在大庭广众之下通过实际事件被传扬。那是怎样的言论呢？正是我常说的那番话。风吹过，叶子飘落了。\BibleQuote{草已枯干，花已凋谢。}\BibleRef{依 40:8}黑夜已逝，白昼已临；虚影被证明是虚妄的，真理显明了。他们升到天上，又降在尘世：因为那汹涌的波涛通过平凡的事件被平息了。怎么做到的呢？所发生的事是一堂课。我对自己说：后人会学会自制吗？还是不出两日，这些事就被遗忘在脑后呢？警告还在他们耳边回响。让我再次宣讲，我还要再说。有什么益处呢？肯定有益处。因为若所有人都不听，一半人会听；若一半人不听，三分之一人会听；若三分之一不听，四分之一；若四分之一不听，也许十人；若十人不听，也许五人；若五人不听，也许一人；若一人也不听，我自己也预备了赏报。\BibleQuote{草必枯干，花必凋谢；但天主的圣言永远常存。}\BibleRef{依 40:8}\par

3. 你们看见世事的虚空吗？你们看见权势的脆弱吗？你们看见那财富吗？我一向称它为逃跑者，不仅逃跑者，而且也是杀人者。因为它不仅遗弃拥有它的人，而且还杀害他们；因为当有人向它献媚时，它最背叛他。你为何向那财富献媚，它今天为你所有，明天却为他人所有？你为何向那永远不能持守的财富献媚？你想向它献媚吗？你想持守它吗？不要埋葬它，而要把它交在穷人手中。因为财富是一只野兽：若紧紧抓住，它就逃跑；若松开，它便留在原地。\Quote{因为，}据说，\BibleQuote{他慷慨施舍，救济穷人；他的义德永远常存。}那么，分散它，好使它与你同在；不要埋葬它，免得它逃跑。财富在哪里？我很乐意问那些已离世的人。我现在说这些话，并非指责——断乎不可！——也不是刺痛旧伤，而是努力为你们从他人的沉船中寻得一个避风港。当兵士和刀剑威胁时，当全城怒火沸腾时，当皇帝的威严无能为力、紫袍受辱时，当各处充满狂乱时，那时财富在哪里？你们的银盘在哪里？你们的银榻在哪里？你们的家奴在哪里？他们都已逃走了；阉人在哪里？他们都跑掉了；你们的朋友在哪里？他们变了面孔；你们的房屋在哪里？它们被关闭了；你们的金钱在哪里？它的主人逃跑了；那金钱本身，又在哪里？它被埋藏了。它被藏在哪里？我不断地宣称财富背叛那些滥用它的人，我是否对你们显得苛刻而烦扰？现在时机已到，证明我话的真实。你们为何如此紧紧抓住它，在考验时它却于你们无益？如果它在你陷入困境时有力量，就让它在困境中援助你；但如果它在那时逃跑，你们需要它做什么？事件本身作了见证。它有何益处？刀剑已磨利，死亡临近，军队咆哮：对迫在眉睫的危险充满恐惧；但财富却无处可见。那逃跑者逃往何处？它本是这一切邪恶的根源，但在急难时刻却逃跑了。然而许多人指责我，不断说：你总是指责富人；而他们却指责穷人。不错，我确实指责富人；或者更确切地说，不是富人，而是那些滥用财富的人。因为我不断地说，我并不攻击富人的品性，而是攻击贪婪之人的品性。富人是另一回事，贪婪之人是另一回事；丰裕之人和贪吝之人是不同的。要清楚区分，不要混淆不同的事物。你是富人吗？我不禁止你。你是贪婪之人吗？我谴责你。你有自己的财产吗？享用它。你夺取他人的财产吗？我绝不缄默。为此你们要用石头砸我吗？我准备流我的血；只禁止你们的罪。我不在意仇恨，不在意战争；我只在意一件事：听众的进步。富人是我的孩子，穷人也是我的孩子；同一个母腹孕育了这两者，两者是同一阵痛所生的。那么，如果你指责穷人，我就谴责你；因为穷人受的损失不如富人的大。因为对穷人没有多大损害，他所受的伤害仅限于金钱；但对你而言，伤害触及灵魂。谁愿意弃绝我，就让他弃绝吧；谁愿意用石头砸我，就让他砸吧；谁愿意恨我，就让他恨吧；因为敌人的阴谋为我乃是胜利冠冕的保证，我的赏报将如我的伤痕一样众多。\par

4. 所以我不怕敌人的阴谋：我只怕一件事，那就是罪。若无人指控我有罪，那么就让全世界与我为敌吧。因为这种战争只会使我更加昌盛。我也愿以此教导你们。不要畏惧权贵的计谋，而要畏惧罪的能力。若你尚未伤害自己，就没有人能伤害你。若你没有罪，即使千万把刀剑威胁你，天主也会将你从他们手中救出；但若你有罪，即使你在乐园里，也会被赶出去。亚当在乐园里，却堕落了；约伯在粪堆上，却得胜加冕。乐园对前者有何益处？粪堆对后者有何损害？无人给前者设陷阱，他却倾覆了；魔鬼给后者设陷阱，他却得胜了。魔鬼不是夺走了他的财产吗？是的，但没有夺走他的虔敬。不是残害了他的儿子吗？是的，但没有动摇他的信德。不是撕裂了他的身体吗？是的，但没有找到他的宝藏。不是利用他的妻子攻击他吗？是的，但没有击倒这位战士。不是向他投掷箭矛吗？是的，但他没有受伤。魔鬼推进了攻城器械，却不能动摇这塔楼；卷起了巨浪，却不能击沉这船。我恳求你们遵守这条法则，是的，我抱住你们的膝盖——若不是用身体的手，也是用精神的手——并流下恳求的眼泪。我恳求你们遵守这条法则，这样就没有人能伤害你们。绝不要称富人为有福的；绝不要称任何人为可怜的，除非他活在罪中；要称那活在义中的人为有福的。因为使人幸福或不幸的，不是环境本身，而是人的心志。如果你的良心不控告你，就绝不要怕刀剑；如果你的良心清白，就绝不要怕战争。那些离去的人在哪里？告诉我。从前不是所有人都向他们俯首吗？那些掌权者不是在他们面前战栗不已吗？不是奉承他们吗？但罪来了，一切都显出了真相；原先的侍从成了审判官，谄媚者变成了刽子手；那些曾亲吻他手的人，亲自把他从教堂里拖出来；昨天吻他手的人，今天成了他的敌人。为什么？因为他昨天也不是真心爱他。时机一到，演员们便露出了真面目。你们昨天不是吻他的手，称他为救主、护卫者和恩人吗？不是无止境地为他写颂词吗？为什么今天又控告他？为什么昨天是赞美者，今天成了告发者？为什么昨天是颂词，今天却是辱骂？这是什么改变？这是什么转变？\par

5. 但我不像这样：他曾阴谋对付我，我却成了他的保护者。我在他手中受了无数苦，却没有报复。因为我效法我老师的榜样，祂在十字架上说：\Quote{宽赦他们罢，因为他们不知道他们做的是什么。}我现在说这些，是希望你们不被恶人的猜疑所败坏。自从我管理这城以来，发生了许多变化，却没有人学会自制？但我说\Quote{没有人}时，并非谴责所有人，天主禁止。因为这片肥沃的土地既然接受了种子，不可能不结出一穗麦子；但我不满足于许多人得救，而是希望所有人都得救。哪怕只有一个人处于灭亡状态，我也灭亡，并且认为牧者应当效法，祂有九十九只羊，仍去追赶那只迷失的。\BibleRef{路 15:4}财富能存到几时？这银子和金子能存到几时？这些酒宴能到几时？奴隶的谄媚能到几时？这些饰有花环的酒杯能到几时？这些撒旦的饮宴，充满魔鬼的作为，能到几时？\par

难道你们不知道，现世的生活是寄居他乡吗？你是公民吗？不，你是过客。你明白我在说什么吗？你不是公民，而是过客和旅人。不要说：我有这个城、那个城。没有人有城。城在天上。现世的生活只是一段旅程。我们每天都在旅行，而本性也在运行。有些人在路上囤积货物，有些人在路上埋藏珠宝。当你进入客栈时，你会装饰客栈吗？不，你只是吃喝，然后赶紧离开。现世的生活就是一家客栈：我们进入了它，我们要结束这现世的生活；让我们渴望带着美好的希望离开，什么也不留在这里，以免在那里失去它们。你进入客栈时，对仆人说些什么？\Quote{小心收好我们的东西，不要留下任何东西在这里，免得丢失，哪怕是微小琐碎的东西，也要全部带回家里。}你是一个过客和旅人，而且确实比过客更微不足道。怎么讲？我来告诉你。过客知道自己什么时候进客栈，什么时候离开；因为离开和返回都在他掌控之中；但我进入客栈，即这现世的生活时，却不知道什么时候离开；我可能为自己准备长久的生计，而主人突然召唤我说：\Quote{糊涂人哪，你所备置的，将归谁呢？因为今夜就要索回你的灵魂。}你离开的时间不确定，你的财产保有不安全，你周围有无数的悬崖深渊和惊涛骇浪。你为什么为影子而疯狂？为什么舍弃现实而追逐影子？\par

6. 我说这些话，并将不断说下去，不断造成痛苦，包扎伤口；这不是为了跌倒的人，而是为了那些仍然站立的人。因为他们已经离去，他们的路程结束了，而那些尚站立的人，因他们的灾祸获得了更稳固的地位。你说：\Quote{那么，我们该做什么？}只做一件事：憎恨财富，热爱你的生命——抛弃你的财物；我不是说全部，而是砍掉多余的部分。不要贪图别人的财物，不要剥削寡妇，不要掠夺孤儿，不要强占他的房屋：我不是对个人发言，而是针对事实。但如果有人良心责备他，他自己要负责，不是我的话。你为什么在招致怨恨之处抓取？要在能获得冠冕之处抓取。努力抓住的不是地，而是天。\BibleQuote{天国是以猛力夺取的，以猛力夺取的人就得到了它。}你为什么抓住责备你的穷人？抓住为此赞美你的基督。你看到你的愚昧和疯狂了吗？你抓住那有一点点东西的穷人？基督说：\Quote{抓住我；我为此感谢你，抓住我的王国，用暴力夺取它。}如果你想要抓住地上的王国，或者说如果你对它图谋不轨，你会受罚；但对于天国，如果你不抓住它，也会受罚。涉及世俗事物时，有恶意；涉及属灵事物时，有爱。每天默想这些事，如果两天后你看到另一个人乘着战车，穿着丝绸衣服，骄傲自满，不要再惊慌烦恼。不要赞美富人，只赞美活在正义中的人。不要辱骂穷人，但要学习对一切事有正直准确的判断。\par

不要远离教会，因为没有什么比教会更强。教会是你的希望，你的救恩，你的避难所。它高于天，宽于地。它永不衰老，常保活力。为此，以象征其坚固稳定，\WorkTitle{圣经}称它为山；以象征其纯洁，称它为童贞女；以象征其庄严，称它为王后；以象征其与天主的关系，称它为女儿；为表达其多产，称那生了七个孩子的为荒胎者；事实上，它用了无数名称来代表其高贵。因为教会的元首有许多名字：被称为圣父，以及道路\BibleRef{若 14:6}、生命\BibleRef{若 14:6}、光、手臂、赎罪祭\BibleRef{若一 2:2}、根基\BibleRef{格前 3:11}、门\BibleRef{若 10:7}、无罪的、宝藏、主、天主、子、独生子、天主的形状\BibleRef{斐 2:6}、天主的肖像\BibleRef{哥 1:15}。教会本身也是如此：一个名字足以呈现全部真理吗？绝不。但正因如此，有无数的名字，好让我们能学到一些关于天主的知识，尽管只是很小的一部分。同样，教会也被称为许多名字。她被称为童贞女，虽然她从前是妓女；因为这是新郎所行的奇迹，他娶了妓女，把她变成了童贞女。哦！何等新奇的事！对我们来说，婚姻破坏童贞，但对天主来说，婚姻恢复了童贞。对我们来说，一个童贞女结婚后就不再是童贞女；但对基督来说，一个妓女结婚后就变成了童贞女。\par

7. 让那好奇探究天上生育本质的异端者，说\Quote{圣父如何生了圣子？}，解释这一事实：问他，教会本是妓女，如何成为童贞女？她生了孩子后如何仍是童贞女？保禄说：\BibleQuote{因为我是以天主的妒爱妒爱你们，因为我曾把你们许配给一个丈夫，把你们当作纯洁的童女献给基督。}\BibleRef{格后 11:2}何等的智慧和见识！\BibleQuote{我是以天主的妒爱妒爱你们。}这是什么意思？他说：\Quote{我妒爱你们}，你一个属神的人，也会妒爱吗？他说，我如同天主一样妒爱。天主也有妒爱吗？是的，不是情欲的妒爱，而是爱和热切的妒爱。我是以天主的妒爱妒爱你们。我要告诉你他如何显示他的妒爱吗？他看到世界被魔鬼败坏，就交付自己的儿子来拯救它。因为论及天主的话，与论及我们自己的话，力量不同：例如我们说天主妒爱，天主发怒，天主后悔，天主憎恨。这些词是人的，但具有与天主本性相称的含义。天主如何妒爱？\BibleQuote{我是以天主的妒爱妒爱你们。}\BibleRef{格后 11:2}天主发怒吗？\BibleQuote{上主，求你不要在震怒中责备我。}天主沉睡吗？\BibleQuote{醒来，上主，你为什么沉睡？}天主后悔吗？\BibleQuote{我后悔造了人。}\BibleRef{创 6:7}天主憎恨吗？\BibleQuote{我的灵魂恨恶你们的庆节和月朔。}\BibleRef{依 1:14}哦，不要考虑这些表达的贫乏，而要把握它们的神圣含义。天主妒爱，因为祂爱；天主发怒，不是因为顺从情欲，而是为了惩罚和责打；天主沉睡，不是真的打盹，而是长久忍耐。选择这样的表达。因此，当你听到天主生了圣子时，不要想到分裂，而想到本体的合一。因为天主从我们这里采取了这些词语，我们也从祂那里借用了其他词语，好使我们借此得到尊荣。\par

8. 你们明白我所说的话吗？请仔细听，我所爱的人。有天主的名号，也有人的名号。\newline{}天主从我这里领受了名号，祂也将名号赐给我。祂说：把你的给我，把我的给你。你需要我的名号，我不需要你的，但你需要我的，因为我的本性是纯洁无瑕的，而你是一个有肉身的人，也寻求有形的称谓，好藉着借用你熟悉的表达，你这个被肉身包围的人，能够思想超越你理解的思想。\newline{}祂从我这领受了什么名号，又赐给了我什么名号？祂自己是天主，祂称我为神；祂那边是本质实际，我这只是名号的尊荣：\BibleQuote{我曾说过：你们是神，你们都是至高者的子女。}\BibleRef{咏 82:6}\newline{}这里有的是言语，而彼处却是实在本身。祂称我为神，因这名号我得了尊荣。祂自己被称为人，被称为人子，被称为道路、门、磐石。这些词语是祂从我借用的；另一些则是祂从自己赐予我的。\newline{}为何祂被称为道路？为使你明白，藉着祂我们得以亲近父。\newline{}为何祂被称为磐石？为使你明白信德的安全与坚不可摧。\newline{}为何祂被称为根基？为使你明白祂支撑万有。\newline{}为何祂被称为根？为使你明白在祂内我们得成长的力量。\newline{}为何祂被称为牧人？因为祂喂养我们。\newline{}为何祂被称为羊？因为祂为我们牺牲，成了赎罪祭。\newline{}为何祂被称为生命？因为祂使我们从死中复活。\newline{}为何祂被称为光？因为祂解救我们脱离黑暗。\newline{}为何祂被称为臂膀？因为祂与父同一性体。\newline{}为何祂被称为圣言？因为祂由父所生。\newline{}因为正如我的言语是我精神的后裔，同样，圣子是由父所生。\newline{}为何祂被称为我们的衣服？因为我受洗时穿上了祂。\newline{}为何祂被称为筵席？因为我分享奥迹时以祂为食。\newline{}为何祂被称为房屋？因为我住在祂内。\newline{}为何祂被称为住客？因为我们成了祂的圣殿。\newline{}为何祂被称为头？因为我成了祂的肢体。\newline{}为何祂被称为新郎？因为祂娶了我为新娘。\newline{}为何祂被称为无玷？因为祂以童贞女娶了我。\newline{}为何祂被称为主人？因为我是祂的婢女。\par

9. 请观察教会，正如我所说，她有时是新娘，有时是女儿，有时是童贞女，有时是婢女，有时是王后，有时是不育的妇人，有时是山，有时是园子，有时多结果实，有时是百合花，有时是泉源：她是一切。\newline{}所以，听到这些，我求你们不要以为它们是物质的，而要伸展你们的思绪，因为这类事不可能是物质的。例如：山不是婢女，婢女不是新娘，王后不是婢女，然而教会却是这一切。为什么？因为她们存在的要素不是物质的，而是精神的。因为在物质的领域，这些事物被局限在狭小范围内；但在精神的领域，它们有广阔的作用范围。\newline{}\BibleQuote{王后站在你的右边。}\BibleRef{咏 45:10}\newline{}王后？那位被践踏、贫穷的人如何成了王后？她上升到何处？王后自己高高站在君王身旁。如何？因为君王成了仆人；他本性并非如此，但他成了这样。\newline{}因此，要理解属于天主性的事，辨别属于经世的事。理解祂\Emph{原是}什么，以及祂为了你\Emph{成了}什么，不要混淆不同的事物，也不要把祂的慈爱论证当作亵渎的借口。\newline{}祂是崇高的，她是卑微的：崇高不在于地位，而在于本性。祂的本质是纯洁不朽的：祂的本性是不朽的、不可理解的、不可见的、不可领悟的、永恒的、不变的，超越天使的本性，高于天上的能力，压倒理性，超越思想，不是凭眼见而是唯独凭信德领悟的。天使看见祂就战栗，革鲁宾敬畏地用翅膀遮住自己。\newline{}祂俯视大地，大地就震动；祂斥责大海，使之干涸：\BibleRef{依 51:10}祂使河流从旷野涌出；祂用秤称群山，用天平称山谷。\newline{}我该如何表达？我该如何呈现真理？祂的伟大无边，祂的智慧不可估量，祂的判断不可追迹，祂的道路不可探究。这就是祂的伟大和能力，若说使用这些措辞是安全的话。\newline{}但我该怎么办呢？我是人，我用人的言语说话：我的舌头是属土的，我恳求主的宽恕。我并非以傲慢的心使用这些措辞，而是因为我的资源贫乏，源自我的软弱和我们人言的本质。主啊，求你仁慈待我，因为我说这些话并非傲慢，而是因为我无其他言辞可说；尽管如此，我不满足于言语的卑贱，而是乘我理解的翅膀向上翱翔。\newline{}这就是祂的伟大和能力。我说这些，是为了使你本人也学会以同样的方式行动，而不停留在言辞上或措辞的贫乏上。我这样做，你为何惊奇，既然当祂想要向我们呈现超越人类能力的事物时，祂也这样做？因为祂对世人说话，也使用人的比喻，这些比喻确实不足以表达所说的事，也不能完全展现事物的原貌，但对听众的软弱而言已足够。\par

10． 你们要努力，不要因我冗长的讲论而厌倦。因为当祂显示自己时，祂并不像祂真实的那样显示，祂的纯粹本质并未显示（因为没有人见过天主的本性；因为当祂仅被部分启示时，革鲁宾就颤抖——山岳冒烟，海洋干涸，天空震动，如果启示不是部分的，谁能承受呢？）正如我所说的，祂并不显示祂真实的本相，而只按观看者所能看见的程度显示，因此祂有时显现为老年人的形状，有时为少年，有时在火中，有时在空气中，有时在水中，有时披戴盔甲，并非改变祂的本质，而是使祂的显现适应那些受影响者的各种状况。同样，当任何人想要谈论祂时，他也使用人的比喻。例如我说：\BibleQuote{他上了山，在他们面前变了容貌，他的面貌发光如同太阳，他的衣服洁白如雪。}据说，他显示了少许的天主性，他向他们显示了住在他们中间的天主，\BibleQuote{并在他们面前变了容貌。}请仔细注意这句话。作者说：\BibleQuote{他在他们面前变了容貌，他的衣服发光如同光，他的面貌如同太阳。}当我说\Quote{祂的伟大和能力是如此}并加上\Quote{主，求你怜悯我}（因为我对这表达并不满意，反而困惑，因为没有其他现成的措辞）我希望你们明白，我是从圣经中学到这教训的。于是福音作者想要描述他的光辉，他说：\BibleQuote{他发光了。}他如何发光？告诉我。极其明亮。你如何表达这？他发光\BibleQuote{如同太阳。}你说如太阳？是的。为什么？因为我不知道任何其他更明亮的光体。你又说他又洁白如雪？为什么如雪？因为我不知道任何其他更白的物质。因为他并非真的那样发光，这由随后的事证实：门徒们倒在地上。如果他发光如同太阳，门徒们不会跌倒；因为他们每天都看见太阳，并没有跌倒；但既然他发光比太阳或雪更加灿烂，他们因无法承受这光辉而倒在地上。\par

11． 那么，请告诉我，福音作者，他发光比太阳更明亮，而你们却说\BibleQuote{如同太阳？}是的：为了让你认识那光，我不知道任何其他更大的光体，我没有其他在光体中占据君王地位的比较。我说这些，为使你们不要满足于所用语言的贫乏：我已向你们指出门徒们的跌倒：他们倒在地上，惊呆并沉睡。\BibleQuote{起来}他说，并扶起他们，但他们仍被困。因为他们不能忍受那照耀的过度明亮，沉重的睡意占据了他们的眼睛：那显示的光如此远超太阳的光。然而福音作者却说\BibleQuote{如同太阳，}因为那光体为我们所熟悉，并超越所有其余。\par

但正如我所说，那位如此伟大和全能的主，竟渴望一个妓女。我以此名指称我们的人性。如果一个人渴望妓女，他会被定罪，那么天主也渴望妓女吗？是的，确实。再者，人渴望妓女是为了成为奸淫者；但天主渴望妓女，是为了将妓女转变为贞女：因此，人的渴望是所渴望者的毁灭，但天主的渴望是所渴望者的救恩。那位如此伟大和全能的主为何渴望一个妓女？为能成为她的丈夫。祂如何行动？祂不派遣任何仆人到她那里，不派遣天使、总领天使、革鲁宾或色辣芬；而是亲自接近那爱她的。再者，当你听到爱时，不要以为它是感性的。从话语中提取所含的思想，如同优秀的蜜蜂停在花朵上，采蜜而留下药草。天主渴望了一个妓女，祂如何行动？祂不将她引向高处；因为祂不会将妓女带入天堂，而是亲自降下。既然她不能升到高处，祂就降到地上。祂来到妓女那里，并不羞愧；祂来到她隐秘的居所。祂看见她醉酒的状态。祂如何来呢？不是以祂原始本性的纯粹本质，而是成为妓女所是的——并非意图上，而是实际上成为这样，好使她见到祂时不害怕，不会逃走。祂来到妓女那里，成为人。祂如何成为人？祂在母腹中受孕，渐渐长大，像我一样经历人类的成长过程。谁做了这事？是显明的神性，而非天主性本体；是仆人的形象，而非主人的形象；是属于我的肉体，而非属于祂的本质：祂渐渐长大，与人类交往。虽然祂发现妓女——人性——满身疮痍、兽性化、被魔鬼压制，祂如何行动？祂走近她。她看见祂就逃跑。祂召唤贤士们说：\Quote{你们为什么害怕？我不是审判官，而是医生。\BibleQuote{我来不是为审判世界，而是为拯救世界}} \BibleRef{若 12:47}祂立刻召唤贤士们。哦！新奇而奇异的事件。祂来临的最初成果是贤士们。那维持世界者卧在马槽中，那照顾万有者是襁褓中的婴儿。圣殿被建立，天主居住其中。贤士们前来，立即朝拜祂；税吏前来，被转变为福音传道者；妓女前来，被转变为童贞女；客纳罕妇人来，分享了祂的慈爱。这是爱者的标志：不追究罪过，赦免过犯和冒犯。祂如何行动？祂接纳罪人，与她结合。祂给她什么？一枚戒指。什么样的？圣神。保禄说：\BibleQuote{那坚定我们与你们一同归于基督的，就是天主，祂也印封了我们，并赐下圣神作为凭据。} \BibleRef{格后 1:21}于是祂赐予圣神。然后祂说：\Quote{我岂没有将你栽种在园中吗？}她说：\Quote{是。}那么你怎么从那里堕落了？\Quote{魔鬼来把我赶出了园子。}你被栽种在园中，他把你赶出；看，我将你栽种在我内，我扶持你。如何？魔鬼不敢靠近我。我也不是将你带到天上；但这里比天大：我将你带在我内，我是天堂的主。牧人背负你，狼不再来；或者，我容许他靠近。因此主背负我们的人性：魔鬼靠近并被打败。\Quote{我已将你栽种在我内：}因此祂说\BibleQuote{我是葡萄树，你们是枝条：}所以祂将她栽种在祂内。\Quote{但是，}她说，\Quote{我是罪人，不洁。}\Quote{不要让这困扰你，我是医生。我知道我的器皿，我知道它如何被扭曲。它原是泥土器皿，被扭曲了。我藉重生的洗涤重新塑造它，并把它放在火中锻炼。}因为请看：祂从地上取尘土造人；祂塑造了他。魔鬼来，扭曲了他。然后主来，又取了他，重新塑造、重铸祂在圣洗中，祂不让他的身体仍是泥土，而是使之成为更坚硬的器皿。祂将软泥置于圣神的火中。\BibleQuote{祂要以圣神和火洗你们：} \BibleRef{玛 3:11}祂以水受圣洗，为能重新塑造；以火受圣洗，为能坚硬。因此先知受天主引导预先宣告说：\BibleQuote{你必将他们打碎，如同陶匠的器皿。}祂没有说如同每人所有的瓦器：因为陶匠的器皿是指陶匠在轮上正在制作的器皿：陶匠的器皿是泥土的，但我们的却是更坚硬的器皿。因此预先论及藉圣洗所行的重塑，祂说：\BibleQuote{你必将他们打碎，如同陶匠的器皿}——祂的意思是祂重新塑造并重铸它们。我下到圣洗的水中，我天性的样式被重新塑造，圣神的火重铸它，它变成更坚硬的器皿。我的话并非虚夸，且听约伯说：\BibleQuote{祂使我们如同泥土，} \BibleRef{约 10:9}保禄也说：\BibleQuote{但我们有这宝贝在瓦器中。} \BibleRef{格后 4:7}但请思考这瓦器的力量：因为它是被圣神，而不是被火坚硬的。如何证明它是瓦器？\BibleQuote{五次受鞭打四十下减一下，三次受杖击，一次被石头砸，}而这瓦器并未破碎。\BibleQuote{一昼一夜我在深渊中。}他在深渊中，瓦器却没被溶解；他遭遇船难，宝贝却没有失落；船沉了，货物却浮着。\BibleQuote{但我们有这宝贝，}他说。什么样的宝贝？圣神的供应、正义、圣化、救赎。请告诉我，什么样的？\BibleQuote{因耶稣基督之名，起来行走。} \BibleRef{宗 3:6} \BibleQuote{艾尼阿，耶稣基督治好你了，} \BibleRef{宗 9:34} \BibleQuote{我对你说，邪魔，从他身上出去。} \BibleRef{宗 16:18}\par

12. 你曾见过比君王宝藏更璀璨的宝藏吗？因为君王的珍珠所能做的，怎能与宗徒的话语所成就的相比？纵使在死人身上加冕无数，他们也不会复活；但宗徒的一句话发出，便召回了已被取消的本性，使其恢复原状。\Quote{但我们有这宝藏。}哦，宝藏！它不仅被保存，而且保存了存放它的房屋。你明白我所说的吗？地上的君王和统治者拥有宝藏时，会准备宽大的房屋，有坚固的墙壁、门闩、守卫和门锁，好使宝藏得以保存；但基督却反其道而行：祂将宝藏不是放在石器中，而是放在瓦器中。既然宝藏如此伟大，为何容器如此脆弱？但容器脆弱的原因并非因为宝藏伟大；因为这宝藏不是被容器保存，而是它本身保存容器。我存放这宝藏：谁还能偷走它？魔鬼来了，世界来了，群众来了，却没有偷走宝藏；容器被鞭打，宝藏却没有被出卖；它被沉入海中，宝藏却没有沉没；它死了，宝藏却存活。因此祂赐下了圣神的凭据。那些亵渎圣神威严的人在哪里？请留意。\BibleQuote{那在基督内坚固我们和你们，并给我们傅油的，就是天主；祂在我们身上盖了印，并赐下圣神作为凭据。}\BibleRef{格后 1:21}你们都知道，凭据是整个事物的一小部分；让我告诉你们如何理解。有人要出高价买一所房屋，他说：给我一个凭据，好让我有信心；或者有人为自己娶妻，他商谈嫁妆和财产，并说\Quote{给我一个凭据。}请注意：在购买奴隶和一切盟约中都有凭据。既然基督与我们订立了盟约（因为祂要娶我为新娘），祂也给我指定了嫁妆，不是金钱，而是血。但这嫁妆是祂所赐的好东西，\BibleQuote{是眼所未见，耳所未闻，人心所未曾想到的。}\BibleRef{格前 2:9}祂指定这些作为嫁妆：不朽、与天使一同赞美、脱离死亡、摆脱罪恶、继承王国（祂的财富如此巨大）、正义、圣化、脱离现世的邪恶、发现未来的福乐。我的嫁妆何等伟大！现在请专心留意：注意祂做了什么。祂来娶那娼妓——我这样称呼她，虽然她不洁，好让你明白新郎的爱。祂来了；祂接纳了我；祂给我指定了嫁妆：祂说\Quote{我把我的财富给你。}如何给？祂说：\Quote{你失去了乐园？把它取回。你失去了你的美丽？取回它；取回这一切。}然而，这嫁妆并没有在此世给予我。\par

13. 请注意，这就是祂预先谈及这嫁妆的原因；祂在嫁妆中向我保证了身体的复活——不朽。因为不朽并不总是随着复活而来，两者是不同的。因为许多人复活了，却又再次倒下，如拉匝禄和圣人们的身体。但在此并非如此，而是应许了复活、不朽、与天使欢乐相聚的席位、在云中迎接人子，以及\BibleQuote{这样我们就常与主同在}\BibleRef{得前 4:17}这句话的实现，逃脱死亡，摆脱罪恶，彻底摧毁毁灭。那是怎样的一种情形？\BibleQuote{眼所未见，耳所未闻，人心所未曾想到的，就是天主为爱祂的人所预备的。}你赐给我我所不知道的好事？祂说：\Quote{是的；只愿你在此世与我订婚，在这个世界上爱我。}\Quote{那你为何不在此世给我嫁妆？}当你到我父那里，当你进入王宫时，就会给你。你到我这里来！不，是我到你那里去。我来不是要你留在此地，而是要带你回去。不要在此世寻求嫁妆：一切都依靠希望和信德。\Quote{你在这个世界上什么都不给我吗？}祂回答说：\Quote{接受一个凭据，好让你相信我关于将来的事：接受订婚的礼物和聘礼。}因此保禄说：\BibleQuote{我曾把你们许配给一个丈夫。}\BibleRef{格后 11:2}天主将现世的福乐赐给我们作为订婚礼物：它们是未来的凭据；但完全的嫁妆停留在另一个世界。何以见得？我告诉你。在此世我衰老，在那里我不衰老；在此世我死亡，在那里我不死亡；在此世我忧伤，在那里我不忧伤；在此世有贫穷、疾病和阴谋，在那里没有这类事物；在此世有黑暗和光明，在那里只有光明；在此世有阴谋，在那里有自由；在此世有疾病，在那里有健康；在此世有终点的生命，在那里有永不终结的生命；在此世有罪恶，在那里有正义，罪恶被驱逐；在此世有嫉妒，在那里没有这类事物。\Quote{请把这些给我吧！}有人说；\Quote{不！要等待，好让你的同仆也得救；等待，我说。那坚固我们并赐给我们凭据的——}怎样的凭据？圣神，圣神的恩赐。让我谈谈圣神。祂将权戒赐给宗徒们，说\Quote{拿着这个，把它交给众人。}那么这戒指是分给人，却没有分开？正是如此。让我教导你圣神恩赐的意义：伯多禄领受了，保禄也领受了圣神。祂走遍世界，释放罪人脱离罪恶，使跛子复原，给赤身者衣服，复活死人，洁净癞病人，辖制魔鬼，扼杀恶魔，与天主交谈，建立教会，将庙宇夷为平地，推翻祭台，摧毁邪恶，建立德行，使人成为天使。\par

14. 这一切我们曾是。但\Quote{抵押}充满了全世界。当我说全世界时，我指的是太阳所照耀的一切地方：海洋、岛屿、山岳、山谷与丘陵。保禄如同有翅膀的生物般到处奔波，仅凭一张口与仇敌争战；他本是帐棚匠，操持工匠刀，缝合兽皮，然而这手艺并未阻碍他的德行，反倒使帐棚匠比魔鬼更强，无口才者比智者更明智。这从何而来？他领受了抵押，带着那戒指并随身携带。众人都看见君王已聘娶了我们的本性：魔鬼看见便退却，它看见抵押便战栗退缩；它看见宗徒的衣裳\BibleRef{宗 19:11}便逃遁。啊！圣神的能力。祂不仅将权柄赐予灵魂与身体，甚至赐予衣裳；不仅如此，还赐予影子。伯多禄四处行走，他的影子驱散疾病\BibleRef{宗 5:15}，驱逐魔鬼，使死者复活。保禄周游世界，斩除不敬神道的荆棘，广撒虔敬的种子，如同优秀的农夫手持教义的犁头。他向谁而去？向特拉基亚人、西提亚人、印度人、毛里人、撒丁人、哥特人、野蛮人，并将他们全都改变。凭借什么？凭借\Quote{抵押}。他如何足以成就这些事？藉着圣神的恩宠。他无技能、衣不蔽体、赤足，却由那\Quote{也赐给我们圣神作抵押的}所扶持。因此他说：\BibleQuote{这些事谁能当得起呢？}\BibleRef{格后 2:16} \BibleQuote{但我们所能承担的，乃是出于天主，祂使我们能够成为新约的仆役，不是文字，乃是圣神。}\BibleRef{格后 3:5} 请看圣神成就了什么：祂发现大地充满魔鬼，却使其变为天堂。不要思虑眼前之事，当在思想中回顾过去。昔日到处是哀号、祭坛、燔祭的烟雾与气味、污秽的仪式与奥秘、祭献、魔鬼狂欢的场所、魔鬼的堡垒、头戴花冠的淫乱；而保禄独自站立。他如何未被淹没或撕碎？他如何能开口？他进入特拜地区，掳获了人；他进入王宫，使君王成为门徒；他进入审判厅，法官对他说：\BibleQuote{你几乎劝我成为基督徒}，法官便成了门徒；他进入监狱，带走了狱卒；他访问野蛮人岛屿，使毒蛇成为他教训的工具；他访问罗马，使元老院归向他的道理；他访问世界各地的江河与荒野。没有土地或海洋未曾分享他劳苦的益处；因为天主将祂戒指的抵押赐给了人性，当祂赐予时说：\Quote{有些我现今赐给你们，有些我应许。}因此先知论及她说：\BibleQuote{王后佩戴金线绣的衣裳，站在你右边。}祂指的不是真正的衣裳，而是德行。因此圣经在别处说：\BibleQuote{你进来这里，怎么没有穿礼服呢？}所以此处祂指的不是衣裳，而是淫乱、污秽不洁的生活。正如污秽的衣裳象征罪恶，金线绣的衣裳则象征德行。但这衣裳属于君王。祂亲自将衣裳赐予她：因她原是赤身露体、丑陋不堪。\BibleQuote{王后佩戴金线绣的衣裳，站在你右边。}祂说的不是衣裳，而是德行。请注意：这表达本身具有极大的高贵意味。祂不说\BibleQuote{金衣}，而说\BibleQuote{金线绣的衣裳}。请明智聆听。金衣是全然为金的衣裳；但金线绣的衣裳是部分为金、部分为丝的。那么，为何祂说新娘穿的并非金衣，而是金线绣的衣裳呢？请留心。祂指的是教会那以多种形式构成的结构。因为我们并非都属于同一种生活状态，有人是贞女，有人是寡妇，有人过着虔敬的生活——因此教会的袍服象征着教会的构成。\par

15. 既然我们的主知道，如果祂只为我们开辟一条道路，许多人必会退缩，因此祂开辟了各种道路。你或许无法通过童贞之道进入天国，那就通过初婚之道进入吧。你无法通过初婚之道进入？也许你可以通过再婚之道进入。你无法通过节欲之道进入？那就通过施舍之道进入吧：或者你无法通过施舍之道进入？那就试试斋戒之道。如果你不能走这条路，就走那条——或者如果那条不行，就走这条。因此先知所说的不是一件金衣，而是一件金线织成的衣服。它有丝绸、紫色或金色。你无法成为金色部分？那就成为丝绸部分。只要你穿上我的衣服，我就接纳你。因此保禄也说：\Quote{若有人在这根基上建造，用金银、宝石。} \BibleRef{格前 3:12} 你无法成为宝石？那就成为金子。你无法成为金子？那就成为银子，只要你是立在根基上。又在别处说：\Quote{太阳的光荣是一样，月亮的光荣是另一样，星辰的光荣又另一样。} \BibleRef{格前 15:41} 你无法成为太阳？那就成为月亮。你无法成为月亮？那就成为星辰。你无法成为大星？满足于做一颗小星吧，只要你在天上。你无法成为童贞者？那就贞洁地过婚姻生活，只要留在教会内。你无法一无所有？那就施舍，只要留在教会内，只要穿着合适的衣服，只要服从那位王后。这衣服是金线织成的，纹理多样。我决不阻拦你：因为丰富的德行使君王的安排易于实行。\Quote{身穿金线织成的衣服，纹理多样。} 她的衣服是多样的：请展开这个表达中的深意，定睛在这件金线织成的衣服上。因为在这里，有些人过独身生活，另一些人在尊贵的婚姻状态中生活，并不比前者逊色太多：有些人结过一次婚，另一些人在花样年华成为寡妇。乐园有何用途？为何有各样的花卉、树木和许多珍珠？有许多星辰，但只有一个太阳；有许多生活方式，但只有一个乐园；有许多圣殿，但只有一位万殿之母。有身体、眼睛、手指，但这一切组成一个人。在小、大和更小之间也有同样的区别。童贞者需要已婚妇女，因为童贞者也是婚姻的产物，为的是婚姻不被她轻视。童贞是婚姻的根：因此一切事物都相互关联，小的与大的相连，大的与小的相连。\Quote{王后身穿金线织成的衣服，纹理多样，站在你的右边。} 接着是：\Quote{女儿，请听！} 向导对新娘说，你即将从你的家前往新郎的家，新郎的本性远超过你。我是新娘的向导。\Quote{女儿，请听。} 她立刻就成为了妻子？是的：因为这里没有任何物质的事。因为祂娶她为妻，爱她如女儿，照顾她如婢女，保护她如童贞女，围护她如花园，珍爱她如肢体：如头般照顾她，如根般使她成长，如牧人般牧养她，如新郎般迎娶她，如赎罪祭般宽恕她，如羔羊般被祭献，如新郎般保持她的美丽，如丈夫般供应她的需要。有许多含义，为的是我们能分享救恩经济的一小部分。\Quote{女儿，请听} 并观看，看那些既属于婚宴又属于神性的事物。女儿，请听。她起初是魔鬼的女儿，大地的女儿，不配拥有大地，如今却成了君王的女儿。这是爱她的那位所愿望的。因为爱者不调查品性：爱不看丑陋：正因如此，它被称为爱，因为它常对丑陋的人怀有感情。基督也是这样。祂看见一个丑陋的人（因为我不能称她为美丽），便爱上了她，使她变得年轻，没有瑕疵或皱纹。噢，何等的新郎！用恩宠装饰新娘的不优雅！女儿，请听！请听并观看！祂说了两件事：\Quote{请听} 和 \Quote{观看}，这两件事取决于你自己，一件在你的眼睛，另一件在你的耳朵。既然她的嫁妆依赖于听（虽然你们中有些人已经敏锐地觉察到了这一点，但让他们等待那些较迟钝的人：我称赞那些预见了真理的人，也体谅那些只跟随他们的脚步的人），既然嫁妆依赖于听——（听是什么意思？信德：因为 \Quote{信德来自听}，信德与实现和实际经验相对）我之前说过，祂将嫁妆分为两部分，一部分给新娘作为定金，另一部分承诺在未来。祂给了她什么？祂给了她罪的赦免、刑罚的免除、成义、圣化、救赎、主的身体、神圣的灵性筵席、死者的复活。因为宗徒们拥有这一切。因此祂赐下一些部分，并应许其他部分。有些是经验和实现，其他则依赖于望德和信德。现在请听。祂赐予了什么？圣洗圣事和祭献。这些是经验到的。祂应许了什么？复活、身体的不朽、与天使的结合、在天使的喜乐团体中的位置、作为祂国度的公民、无罪的生活、\Quote{眼睛未曾看见，耳朵未曾听见，人心也未曾想到的，天主为爱祂的人所预备的} 美善事物。\par

16. 要知道所说的是什么，免得失去它：我正努力使你们能够理解它。那时新娘的聘礼分为两部分：一部分是现在的事，一部分是将来的事；看得见的事和听得见的事，已给予的事和凭信接受的事，已经验的事和将来要享受的事；属于现世生活的事和复活后的事。前者你们看得见，后者你们听得见。那么请注意他对她所说的话，免得你们以为她只收到了前者，尽管前者是伟大、不可言喻、超乎一切理解的。\Quote{女儿，请听，请看；}听后面的事，看前面的事，免得你们说：\Quote{难道我又要依靠希望，又依靠信德，又依靠将来吗？} 看，现在：我给予一些东西，又许诺另一些东西：后者固然依靠希望，但你们要把前者当作抵押，当作凭据，当作其余部分的证据。我许诺给你一个王国：那么让现在的事成为你信任我的依据。你答应给我一个王国吗？是的。我已经给了你更大的一部分，就是那王国的主，因为\BibleQuote{祂既然没有怜惜自己的儿子，反而为我们众人把祂交出了，岂不也把一切与祂一同赐给我们吗？} \BibleRef{罗 8:32} 你答应给我身体的复活吗？是的；我已经给了你更大的一部分。这更大的一部分是什么？罪过的赦免。为什么这是更大的一部分？因为罪恶产生了死亡。我摧毁了根源，难道不能摧毁果实吗？我干涸了树根，难道不能摧毁产物吗？\Quote{女儿，请听，请看。} 我要看什么？死人复活，癞病人洁净，海被平息，瘫痪者恢复力量，乐园被打开，饼大量涌流，罪过被赦免，跛子跳跃，强盗成为乐园的公民，税吏变成传福音者，妓女变得比处女更端庄。听和看。听前者的事，看这些事。从现有的事接受其他事的证明；关于那些，我已经给了你们凭据，就是比它们更好的东西。\Quote{你这话是什么意思？} 这些是我的。\Quote{女儿，请听，请看。} 这些是我给你的聘礼。那么新娘贡献什么？让我们看看。请问你带来什么，免得你没有嫁妆？我能从异教的祭坛、祭物的烟雾和魔鬼那里带来什么？我能贡献什么？什么？你说？你的意志和你的信德。\Quote{女儿，请听，请看。} 那么你希望我做什么？\Quote{忘记你的民族。} 什么民族？魔鬼，偶像，祭物的烟雾、蒸汽和血。\Quote{忘记你的民族和你的父家。} 离开你的父亲，跟随我。我离开我的父，来到你这里，难道你不愿意离开你的父亲吗？但是当\Quote{离开}这个词用于圣子时，不要理解为实际的离开。它的意思是\Quote{我屈尊俯就，我迁就你们，我取了人的肉身。} 这是新郎和新娘的职责：你应该放弃你的父母，我们应该彼此结合。\Quote{女儿，请听，请看，忘记你的民族和你的父家。} 如果我忘记它们，你给我什么？\Quote{君王就会恋慕你的美丽。} 你有主作为你的爱人。如果你有祂作为你的爱人，你也有属于祂的一切。我希望你们能够理解所说的话：因为这是一个微妙的思想，我想堵住犹太人的口。\par

现在请你们集中精神，我恳求：因为无论人听或不听，我都要挖掘和耕耘土壤。\Quote{女儿，请听，请看，也要忘记你的民族和你的父家，君王就会恋慕你的美丽。} 犹太人从这段经文中理解的是感官的美，不是属灵的美，而是肉体的美。\par

17. 请留意，让我们学习什么是肉身之美，什么是精神之美。有灵魂和身体：它们是两种实体；有身体之美，也有灵魂之美。什么是身体之美？修长的眉毛、愉悦的目光、红润的面颊、鲜红的嘴唇、挺直的颈项、卷曲的长发、纤细的手指、匀称的身材、白皙红润的肤色。这种身体之美是来自本性，还是来自选择？显然是来自本性。请留意，以便你们学习哲学家的概念。这美貌，无论是面容、眼睛、头发、眉毛，是来自本性还是来自选择？很明显是来自本性。因为不优雅的女子，即使千方百计修饰，也无法使身体变得优雅：因为本性的条件是固定的，并被它们无法超越的界限所限制。因此，美丽的女子总是美丽的，即使她对美没有兴趣；而不优雅的女子无法使自己优雅，优雅的女子也无法变得不优雅。为什么呢？因为这些来自本性。好了！你们已经看到了身体之美。现在让我们转向内在，去看灵魂：让婢女靠近主母！让我们转向灵魂。注视那美，或者更确切地说，倾听它：因为你们看不见它，它是无形的——倾听那美。那么什么是灵魂之美？节制、温和、施舍、爱、兄弟般的友爱、柔情、服从天主、遵守法律、正义、忏悔的心。这些是灵魂之美。那么这些不是本性的结果，而是道德倾向的结果。没有这些的人能够获得它们，拥有这些的人如果变得粗心，就会失去它们。因为正如我在身体方面所说的，不优雅的女子无法变得优雅；同样，在灵魂方面，我说相反：不优雅的灵魂可以变得充满恩宠。因为有什么比保禄的灵魂——当他是一个亵渎者和辱骂者时——更不优雅呢？又有什么比他说\Quote{这场好仗，我已打完；这场赛跑，我已跑到终点；这信仰，我已保持了}\BibleRef{弟后 4:7}时更充满恩宠呢？有什么比强盗的灵魂更不优雅呢？又有什么比当他听到\Quote{我实在告诉你，今天你就要与我一同在乐园里}\BibleRef{路 23:43}时更充满恩宠呢？有什么比税吏——当他勒索时——更不优雅呢？但当他表明决心时，又有什么更充满恩宠呢？\BibleRef{路 19:8}你们看，你们无法改变身体的恩宠，因为它不是道德倾向的结果，而是本性的结果。但灵魂的恩宠来自我们自己的道德选择。你们现在已经接受了定义。它们是怎样的？灵魂之美来自对天主的服从。因为如果不优雅的灵魂服从天主，它就脱去不优雅，变得充满恩宠。有人对保禄说：\Quote{扫禄！扫禄！你为什么迫害我？}他回答说：\Quote{主，你是谁？}\Quote{我就是耶稣。}\BibleRef{宗 9:4}他服从了，他的服从使那不优雅的灵魂充满了恩宠。同样，祂对税吏说：\Quote{来跟随我}\BibleRef{玛 9:9}，税吏就起来成了宗徒：那不优雅的灵魂充满了恩宠。从何而来？由于服从。同样，祂对渔夫们说：\Quote{来跟随我，我要使你们成为渔人的渔夫}\BibleRef{玛 4:19}；因着他们的服从，他们的心灵充满了恩宠。那么，让我们看看祂在此处所说的是哪种美。\Quote{女子，请听，请看，请忘记你的民族和你的父家，君王就会恋慕你的美色。}祂将恋慕哪种美？精神的美。为什么？因为她要\Quote{\Emph{忘记}}——祂说，\Quote{请听，请忘记}。这些都是道德选择的行为。祂说：\Quote{请听！}一个不优雅的人听了，而她身体的不优雅并未消除。祂对那个罪妇说\Quote{请听}，如果她服从，她会看到怎样的美赐予了她。既然新娘的不优雅不是身体的，而是道德的（因为她没有服从天主，而是违命），所以祂引领她走向另一种补救。那么，你并非由于本性，而是由于道德选择变得不优雅：你因服从而变得充满恩宠。\Quote{女子，请听，请看，请忘记你的民族和你的父家，君王就会恋慕你的美色。}然后，为了让你们知道祂所指的并非任何感官可见之物——当你们听到美这个词时，不要想到眼睛、鼻子、嘴或脖子，而要想到虔诚、信德、爱德——这些内在的事物，\Quote{因为君王女儿的尊贵全在于内在。}为了这一切，让我们感谢天主的赏赐，因为光荣、尊威、权能都归于祂，直到永远。阿们。\par

\SourceRecord{Translated by W.R.W. Stephens.；Nicene and Post-Nicene Fathers, First Series；Vol. 9.；Edited by Philip Schaff.；Buffalo, NY: Christian Literature Publishing Co.,；1889.；<http://www.newadvent.org/fathers/1915.htm>.}
